\section{Approximation and Generalization Theory}\label{sec:approx}

In this section, we present generic approximation and generalization theory and defer detailed proofs to Section \ref{sec:proof-sketch-approx} and \ref{proof:sketch-generalization} respectively. Firstly, we introduce the approximation theory of utilizing deep neural networks to approximate Hölder functions. The approximation error is determined by effective Minkowski dimension of data distribution and probability of low-density area.  Furthermore, we present the generalization error when  approximating regression function $f^*$. The convergence rate also depends on effective Minkowski dimension. 
%A key property of the network architecture in Definition \ref{defi: nn-class} is its universal approximation ability. \citet{nakada2020adaptive} provided that ReLU neural networks can universally approximate a function in Hölder Space whose support has low intrinsic dimensionality within a proper size. We extend this result to show the approximation ability of these networks are adaptive to data distribution.

%The result is summarized in the following theorem:
%To establish the statistical guarantee of estimating $f^*$, we develop the following neural network approximation theory, which is used for bounding the bias of estimation.
\begin{theorem}[Approximation of deep neural networks]\label{lemma:approx}
    Suppose Assumption \ref{assumption:pdata_cover} hold. 
    For $\beta>0$ and any sufficiently small $\epsilon ,\tau >0$, consider a tuple $(L,B,K)$
    \begin{equation*}
        L=C_1,~ B=O(R_S^{\beta s} \epsilon^{-s}),~ \text{and} ~K=C_2 (R_S d)^p \epsilon^{-p/\beta},
    \end{equation*}
    where $R_S>0$ and $p=p(d^{-1}\epsilon^{1/\beta}/2,\tau)$ are given by Assumption \ref{assumption:pdata_cover},  and
    \begin{align*}
        C_1 = O(d),~
        C_2 = O\big( d^{2+\lfloor \beta \rfloor}\big),~\text{and} ~s=s(\beta).
    \end{align*}
    Then for any $f^* \in\cH(\beta, \cD_{\rm data})$, we have
    \begin{equation*}
        \inf_{f \in \mathcal{F}(L,B,K)} \norm{f-f^*}^2_{L^2(P_{\rm data})} \leq \epsilon^2 + 4\tau.
    \end{equation*}
\end{theorem}

The novelty of Theorem \ref{lemma:approx} is summarized below:

\textbf{Dependence on Effective Minkowski Dimension.} The approximation rate in Theorem \ref{lemma:approx} is $O(K^{-\beta/p})$, which only depends on effective Minkowski dimension $p < d$ and function smoothness $\beta$, but not on ambient dimension $d$. Compared to \citet{YAROTSKY2017103}, our results improves the exponential dependence of neural network size on $d$ to that on $p$. Moreover, unlike \citet{nakada2020adaptive} and \citet{chen2022nonparametric}, our results do not require that data distribution is exactly supported on a low-dimensional structure. Instead, our results can work for data distribution with high-dimensional support as long as its effective Minkowski dimension is relatively small.

\textbf{Relaxation to the $L^2$-error.} The approximation error in Theorem \ref{lemma:approx} is established with respect to the $L^2(P_{\rm data})$ norm, while most of existing works focus on the $L^\infty$ error \citep{YAROTSKY2017103,nakada2020adaptive,chen2019efficient}. Intuitively, it is not necessary for the network class to approximate the function value at each point in the domain $\cD_{\rm data}$ precisely when data distribution is highly concentrated at certain subset. Instead, it suffices to approximate $f^*$ where the probability density is significant, while the error for the low-density region can be easily controlled since the regression function $f^*$ and the neural network class $f \in \cF(L,B,K)$ are bounded.

The benefit of using the $L^2$ error is that, we only need to control the approximation error of $f^*$ within some chosen region $S\subseteq \cD_{\rm data}$. Here $S$ has an effective Minkowski dimension $p$, which ensures that it can be covered by $O(r^{-p})$ hypercubes with side length $r$. Then we design deep neural networks to approximate $f^*$ within each hypercube and thus the network size depends on the number of hypercubes used to cover $S$. This explains why network size in Theorem \ref{lemma:approx} depends on $p$. Meanwhile, the probability out of $S$ is negligible since the data density is low. We further demonstrate that this probability $\tau$ is far less than the approximation error in Section \ref{sec:example}. By this means, we succeed to reduce the network size and at the same time achieve a small $L^2$ approximation error. 

We next establish the generalization result for the estimation of $f^*$ using deep neural networks.
\begin{theorem}[Generalization error of deep neural networks]\label{thm:generalization}
Suppose Assumption \ref{assumption:pdata_cover} holds. Fix any sufficiently small $r,\tau >0$ satisfying $r<R_S$ and $\tau < r^{4\beta}/4$.  Set a tuple $(L,B,K)$ with $C_1,C_2$ and $s$ appearing in Theorem \ref{lemma:approx} as
\begin{equation*}
    L=C_1,~ B=O(R_S^{\beta s} r^{-\beta s}), ~\text{and}~K=C_2 R_S^p r^{-p}
\end{equation*}
with $p=p(r,\tau)$. Let $\hat{f}$ be the global minimizer of empirical loss given in \eqref{eq: obj} with the function class 
$\mathcal{F}=\mathcal{F}(L,B,K)$. Then we have 
\begin{align*}
    \EE \| \hat{f} - f^*\|^2_{L^2(P_{\rm data})} = O\biggl(   \tau +  \sigma r^{2\beta} + \frac{\sigma^2}{n}  \biggl(\frac{R_S}{r} \biggr)^p \log\biggl( \frac{(R_S/r)^p}{ r^{4\beta}-4\tau}\biggr)\biggr),
\end{align*}
where $O(\cdot)$ hides polynomial dependence on $d$.
%\begin{equation*}
%    \| \hat{f} - f^*\|^2_{L^2(P_{\rm data})} = \tilde{O} \bigg( \tau +  r^{2\beta} + \frac{1}{n} r^{-p} \log\bigg( \frac{1}{r}\bigg) + \frac{1}{n}\bigg),
%\end{equation*}
%holds with probability at least $1-2\exp(-n^{p/(2\beta+p)})$ for any $n \geq N$ with a sufficiently large $N$.
\end{theorem}

Theorem \ref{thm:generalization} is a statistical estimation result. It implies that the generalization error also depends on effective Minkowski dimension $p$. To establish this result, we decompose the squared error into a squared bias term and a variance term. The bias is tackled with the approximation error in Theorem \ref{lemma:approx} and the variance depends on the network size. With the network size growing, the variance term increases while the bias term decreases, since the approximation capability of neural networks is enhanced as the size of the network enlarges. Therefore, we need to trade off between the squared bias and the variance to minimize the squared generalization error.

Notably, our analysis in Section \ref{sec:approx} holds for any sufficiently small $\tau$ and $r$, and every pair of $\tau$ and $r$ determines a $p$. As shown in Assumption \ref{assumption:pdata_cover}, if $\tau$ and $r$ decreases, the covering number will become larger while the approximation can be more accurate. In order to establish an explicit bound, we need to trade off $\tau$ and $r$ for the given sample size $n$. Therefore the “optimal” $p$ eventually becomes functions of $n$. We call such an “optimal” $p$ effective Minkowski dimension.

In the next section, we give two specific classes of Gaussian random designs to illustrate how effective Minkowski dimension $p(r,\tau)$ scales with $r$ and $\tau$. We further show that, under a proper selection of the region $S$ and the covering accuracy $r$, the convergence rate for the estimation of $f^*$ using deep neural networks is $\tilde{O}(n^{-2\beta/(2\beta+p)})$, where the effective Minkowski dimension $p$ is properly chosen. 
